\section{Гауссовская аппроксимация RR+PR}\label{sec:theory}

Для достаточно малого $w$ совместная цепь
$(\theta_k^{(w)},Z_{k+1})$ имеет единственный инвариантный закон
\cite[Theorem~1]{Levin2026}. Обозначим через $\theta_\infty^{(w)}$
его $\theta$-координату. Из разложения
\[
\E\theta_\infty^{(w)}
=\theta^\star+w\Delta+O(w^{3/2})
\]
следует, что комбинация \eqref{eq:rr} сокращает член $w\Delta$
\cite[Corollary~1]{Levin2026}. Для распределительного результата одного
сокращения ожиданий недостаточно: требуется отделить ведущую марковскую сумму
от нелинейного остатка и сравнить её дисперсию с \eqref{eq:sigmainfty}.

\subsection{Веса и детерминированная ковариационная прокси}

Положим $B_w=I-w\Abar$. Вклад шума, возникшего в момент $l$, в среднее по
окну $k=n_0,\ldots,n-1$ имеет детерминированный вес
\begin{equation}\label{eq:burned-weight}
Q_{l;n_0,n}^{(w)}
=\begin{cases}
\Abar^{-1}(B_w^{\,n_0-l}-B_w^{\,n-l}),&1\leq l<n_0,\\[2mm]
\Abar^{-1}(I-B_w^{\,n-l}),&n_0\leq l\leq n-1.
\end{cases}
\end{equation}
RR-вес равен
\[
Q_l^{\RR}=2Q_{l;n_0,n}^{(\alpha)}
-Q_{l;n_0,n}^{(2\alpha)}.
\]
Из \eqref{eq:contraction} и геометрических сумм следуют фиксированные
константы\linebreak $C_Q,C_V<\infty$, не зависящие от $n$, $n_0$ и $\alpha$, для которых
\begin{equation}\label{eq:weight-controls}
\max_{1\leq l<n}\norm{Q_l^{\RR}}\leq C_Q,\qquad
\sum_{l=1}^{n-2}\norm{Q_{l+1}^{\RR}-Q_l^{\RR}}
\leq C_V/a^2.
\end{equation}
Поэтому ведущая, линейная по $\eps$ часть нормированной ошибки имеет вид
\begin{equation}\label{eq:leading-weighted-sum}
W_{n,n_0}^{\RR}
=-\frac1{\sqrt m}\sum_{l=1}^{n-1}Q_l^{\RR}\eps(Z_l).
\end{equation}
Сумма содержит и индексы до $n_0$: соответствующие веса экспоненциально
затухают, но не равны нулю. Множитель $1/m$ в следующей ковариационной прокси
появляется из нормировки $1/\sqrt m$ в
\eqref{eq:leading-weighted-sum}, то есть из длины окна PR-усреднения, а не из
числа индексов в шумовой сумме.

\begin{lemma}[Сравнение ковариационной прокси]
\label{lem:variance}
Пусть выполнено \eqref{eq:contraction}, $0<\alpha$,
$2\alpha\leq\alpha_\infty$ и $\norm{\SigmaM}<\infty$. Определим
\begin{equation}\label{eq:sigma-proxy}
\Sigma_{n,n_0}^{\mathrm{bRR}}
=\frac1m\sum_{l=2}^{n-1}
Q_l^{\RR}\SigmaM(Q_l^{\RR})^\top,\qquad
\sigma_{n,n_0}^2(u)
=u^\top\Sigma_{n,n_0}^{\mathrm{bRR}}u.
\end{equation}
Положим $\sigma_{n,n_0}(u)=\sqrt{\sigma_{n,n_0}^2(u)}$.
Тогда существует фиксированная константа $C_{\rm var}<\infty$ такая, что
\begin{equation}\label{eq:variance-bound}
\norm{\Sigma_{n,n_0}^{\mathrm{bRR}}-\Sigma_\infty}
\leq\frac{C_{\rm var}}{m\alpha a},\qquad
|\sigma_{n,n_0}^2(u)-\sigma^2(u)|
\leq\frac{C_{\rm var}\norm u^2}{m\alpha a}.
\end{equation}
Константа зависит только от $Q,\Abar$ и $\norm{\SigmaM}$.
\end{lemma}

\begin{proof}
Пусть $\rho_\alpha=(1-\alpha a)^{1/2}$. Из
\eqref{eq:burned-weight} и \eqref{eq:contraction} получаем
\[
\norm{Q_l^{\RR}-\Abar^{-1}}
\lesssim\rho_\alpha^{\,n-l}\quad(l\geq n_0),\qquad
\norm{Q_l^{\RR}}
\lesssim\rho_\alpha^{\,n_0-l}\quad(l<n_0).
\]
Геометрические суммы первых и вторых степеней этих оценок имеют порядок
$(\alpha a)^{-1}$. В части после разогрева раскрываем
\[
Q_l^{\RR}\SigmaM(Q_l^{\RR})^\top-\Sigma_\infty
\]
по степеням $Q_l^{\RR}-\Abar^{-1}$; предшествующую часть оцениваем по энергии
$\sum_{l<n_0}\norm{Q_l^{\RR}}^2$. Не более двух граничных позиций дают
$O(m^{-1})$, которое поглощается правой частью
\eqref{eq:variance-bound}. После деления на $m$ получаем первое неравенство,
а второе следует из $|u^\top Hu|\leq\norm H\norm u^2$.
\end{proof}

\subsection{Мартингал и моментная оценка остатка}

Пусть $\mathsf P$ --- переходное ядро цепи. Уравнение Пуассона имеет
ограниченное решение
\[
\widehat\eps=\sum_{j\geq0}\mathsf P^j\eps,\qquad
\widehat\eps-\mathsf P\widehat\eps=\eps,\qquad
\norm{\widehat\eps}_\infty\lesssim\tmix\norm\eps_\infty .
\]
Кроме того,
\[
\mathcal V_\eps(z)=\mathsf P(\widehat\eps\widehat\eps^\top)(z)
-(\mathsf P\widehat\eps)(z)(\mathsf P\widehat\eps)(z)^\top,
\qquad \E_\pi\mathcal V_\eps(Z)=\SigmaM.
\]
Для $2\leq l\leq n-1$ положим
\[
\Delta M_l^{\mathrm{bRR}}
=Q_l^{\RR}\{\widehat\eps(Z_l)
-\mathsf P\widehat\eps(Z_{l-1})\},\qquad
M_{n,n_0}^{\mathrm{bRR}}
=\sum_{l=2}^{n-1}\Delta M_l^{\mathrm{bRR}} .
\]
Это мартингальные разности. Суммирование Абеля даёт
\begin{equation}\label{eq:poisson-decomp}
W_{n,n_0}^{\RR}
=-\frac{M_{n,n_0}^{\mathrm{bRR}}}{\sqrt m}+D_{2,n,n_0},
\qquad
\norm{D_{2,n,n_0}}_\infty\lesssim\frac{\tmix}{\sqrt m}.
\end{equation}
Здесь
\[
D_{2,n,n_0}=-\frac1{\sqrt m}\left\{
Q_1^{\RR}\widehat\eps(Z_1)
+\sum_{l=1}^{n-2}(Q_{l+1}^{\RR}-Q_l^{\RR})
\mathsf P\widehat\eps(Z_l)\right\};
\]
его оценка следует из \eqref{eq:weight-controls}.
Здесь и далее $\lesssim$ не скрывает зависимости от
$\tmix,n,p,\alpha$; фиксированные величины
$C_A,\kappa_Q,a,\norm{\Abar^{-1}}$ и $\norm\eps_\infty$ считаются
константами задачи.

Предсказуемая квадратичная вариация этого мартингала равна
\begin{equation}\label{eq:bracket}
\langle M\rangle_{n,n_0}
:=\sum_{l=2}^{n-1}\E[\Delta M_l^{\mathrm{bRR}}
(\Delta M_l^{\mathrm{bRR}})^\top\mid\mathcal F_{l-1}]
=\sum_{l=2}^{n-1}Q_l^{\RR}\mathcal V_\eps(Z_{l-1})(Q_l^{\RR})^\top.
\end{equation}

\begin{lemma}[Концентрация случайной скобки]\label{lem:bracket}
При условиях леммы~\ref{lem:variance} и условиях 1, 3 для всех $p\geq2$,
всех $u$ и любого начального распределения цепи
\begin{equation}\label{eq:bracket-bound}
\norm{u^\top\langle M\rangle_{n,n_0}u
-m\sigma_{n,n_0}^2(u)}_{L_p}
\lesssim C_Q^2\norm u^2\norm\eps_\infty^2
\tmix^{5/2}\sqrt{pn}.
\end{equation}
Здесь и в дальнейшем $\lesssim$ не скрывает зависимости от
$\tmix,n,p,\alpha$.
\end{lemma}

\begin{proof}
Для $2\leq l<n$ положим
\[
g_l(z)=u^\top Q_l^{\RR}
\{\mathcal V_\eps(z)-\SigmaM\}(Q_l^{\RR})^\top u.
\]
Из пуассоновского представления
$\norm{\mathcal V_\eps}_\infty
\lesssim\tmix^2\norm\eps_\infty^2$, поэтому
$\norm{g_l}_\infty\lesssim
C_Q^2\norm u^2\tmix^2\norm\eps_\infty^2$ и $\E_\pi g_l=0$.
Левая часть \eqref{eq:bracket-bound} есть норма
$\sum_{l=2}^{n-1}g_l(Z_{l-1})$. Неоднородное марковское неравенство
\cite[Lemma~11]{Levin2026} даёт дополнительный множитель
$\sqrt{p\tmix n}$ и завершает доказательство.
\end{proof}

\begingroup
\emergencystretch=2em
Для точной формулировки обозначим через
$\alpha_{\mathrm{P2}},\alpha_{\mathrm{C6}},
 \alpha_{\mathrm{P8}},\alpha_{\mathrm{P9}},\alpha_{\mathrm{C4}}$
пороги из Proposition~2, Corollary~6, Propositions~8--9 и Corollary~4
расширенной версии \cite{Levin2026} при используемых ниже порядках, а через
$\alpha_{\mathrm{P5}}(r,\tmix)$ --- порог Proposition~5 для момента порядка
$r$. Положим
\endgroup
\[
\begin{aligned}
\alpha_*(p,q,\tmix)=\min\{&\alpha_{\mathrm{P2}},\alpha_{\mathrm{C6}},
 \alpha_{\mathrm{P5}}(p,\tmix),\alpha_{\mathrm{P5}}(2p,\tmix),\\[-1mm]
 &\alpha_{\mathrm{P8}},\alpha_{\mathrm{P9}},\alpha_{\mathrm{C4}}\}.
\end{aligned}
\]
\begin{equation}\label{eq:alpha-adm}
\begin{aligned}
\alpha_{\mathrm{adm}}(p,q,\tmix)
 &=\min\Biggl\{\alpha_\infty,\frac1{2\norm{\Abar}},
 \alpha_*(p,q,\tmix),\alpha_{\rm prod}(p),\\[-1mm]
 &\hspace{25mm}\alpha_{\rm prod}(2p),\frac1{2a},
 \frac1{2a\tmix}\Biggr\}.
\end{aligned}
\end{equation}
Аргументы остальных исходных порогов выбираются ровно при тех порядках
моментов, при которых они используются, и опущены только для читаемости.

Для $w\in\{\alpha,2\alpha\}$ определим компоненты разложения второго порядка.
Положим $J_0^{(\ell,w)}=0$ и $H_0^{(2,w)}=0$, а также
\begin{align*}
J_k^{(0,w)}&=B_wJ_{k-1}^{(0,w)}-w\eps(Z_k),\\
J_k^{(\ell,w)}&=B_wJ_{k-1}^{(\ell,w)}
-w\widetilde A(Z_k)J_{k-1}^{(\ell-1,w)},\quad \ell=1,2,\\
H_k^{(2,w)}&=(I-wA(Z_k))H_{k-1}^{(2,w)}
-w\widetilde A(Z_k)J_{k-1}^{(2,w)}.
\end{align*}
Тогда повторная подстановка в \eqref{eq:lsa} даёт точное равенство
\begin{equation}\label{eq:depth-two}
\begin{aligned}
\theta_k^{(w)}-\theta^\star
 &=\Gamma_{1:k}^{(w)}(\theta_0-\theta^\star)+J_k^{(0,w)}
   +R_{k,\rm fin}^{(w)},\\
R_{k,\rm fin}^{(w)}
 &=J_k^{(1,w)}+J_k^{(2,w)}+H_k^{(2,w)}.
\end{aligned}
\end{equation}

Нам понадобится один локальный переход от стационарных оценок компонент к
детерминированному старту.

\begin{lemma}[Перенос после разогрева]\label{lem:startup}
Пусть $p\geq2$, $q\geq4p$ и
$2\alpha\leq\alpha_{\mathrm{adm}}(p,q,\tmix)$. Для каждой ветви существует
стационарная двусторонняя версия $R_{k,\rm stat}^{(w)}$ и сопряжение с
конечной версией из \eqref{eq:depth-two}, равномерное по закону $Z_0$, такие,
что
\begin{equation}\label{eq:startup-pointwise}
\norm{R_{k,\rm fin}^{(w)}-R_{k,\rm stat}^{(w)}}_{L_p}
\leq A_{\rm st}(p,q,w)
\exp\!\left(-\frac{c_{\rm st}wak}{p}\right),
\end{equation}
где $c_{\rm st}>0$ фиксирована и
\begin{equation}\label{eq:Ast}
A_{\rm st}(p,q,w)=C_{\rm st}(1+d^{1/q})(p^7+p^8/a)
\tmix^5\sqrt{w/a}\,\log^3\!\frac1{wa}.
\end{equation}
Следовательно,
\begin{equation}\label{eq:startup-window}
\begin{aligned}
&\left\|\frac1{\sqrt m}\sum_{k=n_0}^{n-1}
 \begin{aligned}
 \bigl\{&2(R_{k,\rm fin}^{(\alpha)}-R_{k,\rm stat}^{(\alpha)})\\[-1mm]
 &{}-(R_{k,\rm fin}^{(2\alpha)}-R_{k,\rm stat}^{(2\alpha)})\bigr\}
 \end{aligned}\right\|_{L_p}\\
&\qquad\lesssim\frac{pA_{\rm st}(p,q,\alpha)}{\alpha a\sqrt m}
 \exp\!\left(-\frac{c_{\rm st}\alpha an_0}{p}\right).
\end{aligned}
\end{equation}
\end{lemma}

\begin{proof}
Построим стационарную координату $H^{(2,w)}$. На одной двусторонней
стационарной цепи запустим все $J$- и $H$-координаты из нуля в момент $-N$.
Corollary~4 из \cite{Levin2026} даёт инвариантный закон $J$-координат.
Сравнивая запуски из моментов $-N'$ и $-N$, $N'\geq N$, используем после
$-N$ общую траекторию цепи, то есть сопряжение с $T=0$. Покомпонентный расчёт
из доказательства Proposition~5 применим к этому синхронному сопряжению;
начальная стоимость в нём оценивается Lemmas~4, 8 и Proposition~8 при порядках
$p$ и $2p$. Условие $q\geq4p$ обеспечивает допустимость последнего порядка.
В результате, для $-N\leq r\leq0$ и $\ell=0,1,2$,
\[
\norm{\Delta J_r^{(\ell,w)}}_{L_p}
\leq A_J(p,q,w)
\exp\!\left\{-\frac{wa(r+N)}{12p}\right\},
\]
где
\[
A_J(p,q,w)\lesssim(1+d^{1/q})p^7\tmix^5
\sqrt{w/a}\,\log^3\!\frac1{wa}.
\]
Вычитая две конечнопрошлые $H$-рекурсии, получаем
\[
\Delta H_0
=\Gamma_{-N+1:0}^{(w)}H_{-N,N'}^{(2,w)}
-w\sum_{l=-N+1}^{0}\Gamma_{l+1:0}^{(w)}
\widetilde A(Z_l)\Delta J_{l-1}^{(2,w)}.
\]
Условие~4, Proposition~9 из \cite{Levin2026} и геометрическая свёртка дают
\[
\norm{\Delta H_0}_{L_p}
\lesssim\left\{B_H(p,q,w)+\frac paA_J(p,q,w)\right\}e^{-cwaN/p},
\]
где
\[
B_H(p,q,w)\lesssim(1+d^{1/q})p^{7/2}\tmix^{5/2}
 w^{3/2}\log^{3/2}\!\frac1{wa}.
\]
Следовательно, $H_{0,N}^{(2,w)}$ фундаментальна в $L_p$.
Сдвиг конструкции определяет двустороннюю стационарную координату, а оценка
$B_H$ переходит к пределу по лемме Фату. Обе RR-ветви строятся на одной
двусторонней цепи.

Теперь зададим одно сопряжение конечной и стационарной версий, пригодное для
всех оставшихся оценок. Положим $b=\lceil\tmix\rceil$. В моменты, кратные
$b$, максимально сопрягаем конечные точки двух $\mathsf P^b$-блоков;
условные распределения путей внутри блока существуют, поскольку пространство
польское. После первого совпадения используем общие переходы. Получаем липкое
блочно-коадаптированное сопряжение; момент встречи $T$ является stopping time
блочной фильтрации и
\[
\Pr(T>r)\leq C\exp(-cr/\tmix).
\]
Доказательство Proposition~5 использует для $J$-координат только это свойство,
совпадение траекторий после $T$ и равномерные моменты. Повторяя в треугольных
рекурсиях его покомпонентную подстановку и оценку события $\{T>k\}$ по
Гёльдеру, при порядках $p$ и $2p$ получаем для того же сопряжения
\begin{equation}\label{eq:J-same-coupling}
\sum_{\ell=0}^2
\norm{\Delta J_k^{(\ell,w)}}_{L_{2p}}
\leq A_J(2p,q,w)\exp(-cwak/p).
\end{equation}
Таким образом, здесь не переносится абстрактная Wasserstein-оценка на другое
сопряжение: её покомпонентный расчёт повторён для выбранного блочного
сопряжения.

\begingroup
\emergencystretch=2em
После $T$ общая будущая цепь условно по $\mathcal G_T$ имеет исходное ядро.
Поэтому условие~4 распространяется с детерминированных на
$\mathcal G_s$-измеримые векторы поточечным условливанием. В
\eqref{eq:alpha-adm} обеспечено $wa\tmix\leq1/2$. Если
$\sup_s\norm{V_s}_{L_{2p}}\leq B$, то
\par
\endgroup
\[
\norm{V_s\mathbf1_{\{T=s\}}}_{L_p}
\leq B\Pr(T=s)^{1/(2p)}
\leq CB e^{-cs/(p\tmix)}
\leq CB e^{-c'was/p}.
\]
Следовательно,
\begin{align*}
\norm{\Gamma_{T+1:k}^{(w)}V_T\mathbf1_{\{T\leq k\}}}_{L_p}
&\leq C\sum_{s=0}^k e^{-cwa(k-s)/p}
 \norm{V_s\mathbf1_{\{T=s\}}}_{L_p}\\[-1mm]
&\leq C\frac p{wa}B\,e^{-c'wak/p}.
\end{align*}
А для адаптированных $U_l$, удовлетворяющих
$\norm{U_l}_{L_{2p}}\leq B e^{-c_0wal/p}$, та же условная оценка даёт
\[
w\sum_{l=1}^k
\norm{\Gamma_{l+1:k}^{(w)}U_l\mathbf1_{\{T<l\}}}_{L_p}
\leq C\frac paB e^{-c'wak/p}.
\]

На событии $\{T\leq k\}$ имеем
\[
\Delta H_k^{(2,w)}
=\Gamma_{T+1:k}^{(w)}\Delta H_T^{(2,w)}
-w\sum_{l=T+1}^k\Gamma_{l+1:k}^{(w)}
 \widetilde A(Z_l)\Delta J_{l-1}^{(2,w)}.
\]
Первый член контролируется случайным рестартом с $B=B_H(2p,q,w)$. Во втором
берём
$U_l=\mathbf1_{\{T<l\}}\widetilde A(Z_l)\Delta J_{l-1}^{(2,w)}$ и применяем
\eqref{eq:J-same-coupling} и устойчивую свёртку. На дополнении по Гёльдеру
\[
\norm{\Delta H_k^{(2,w)}\mathbf1_{\{T>k\}}}_{L_p}
\leq C B_H(2p,q,w)e^{-ck/(p\tmix)}
\leq C B_H(2p,q,w)e^{-cwak/p}.
\]
Собирая три координаты и используя определения $A_J,B_H$, получаем
\eqref{eq:startup-pointwise} с \eqref{eq:Ast}. Суммирование по окну и
$1-e^{-c\alpha a/p}\gtrsim\alpha a/p$ дают
\eqref{eq:startup-window}.
\end{proof}

Из \eqref{eq:depth-two}, перестановки сумм в $J^{(0,w)}$ и
\eqref{eq:poisson-decomp} получаем точное разложение
\begin{equation}\label{eq:full-decomp}
\sqrt m\,u^\top
(\bar\theta_{n,n_0}^{\RR,\alpha}-\theta^\star)
=-\frac{u^\top M_{n,n_0}^{\mathrm{bRR}}}{\sqrt m}
+\mathcal R_{n,n_0}^{\RR}(u),
\end{equation}
где
\begin{align}\label{eq:composite-remainder}
\mathcal R_{n,n_0}^{\RR}(u)
&=\frac1{\sqrt m}\sum_{k=n_0}^{n-1}u^\top
\{2\Gamma_{1:k}^{(\alpha)}-\Gamma_{1:k}^{(2\alpha)}\}
(\theta_0-\theta^\star)+u^\top D_{2,n,n_0}\notag\\
&\quad+\frac1{\sqrt m}\sum_{k=n_0}^{n-1}u^\top
\{2R_{k,\rm fin}^{(\alpha)}-R_{k,\rm fin}^{(2\alpha)}\}.
\end{align}


\begin{lemma}[Момент остатка]\label{lem:remainder}
Пусть $p\geq2$ и $q\geq4p$. Предположим, что
$2\alpha\leq\alpha_{\mathrm{adm}}(p,q,\tmix)$. Обозначим
\[
L_\alpha=\log\frac1{\alpha a}.
\]
Тогда
\begin{equation}\label{eq:remainder-bound}
\norm{\mathcal R_{n,n_0}^{\RR}(u)}_{L_p}
\lesssim\norm u\,\mathcal B_R(m,n_0,p,q,\alpha),
\end{equation}
где
\begin{align}
\mathcal B_R
&=
\frac{\norm{\theta_0-\theta^\star}}{\alpha a\sqrt m}
\left\{(1-\alpha a)^{n_0/2}
+p\exp\!\left(-\frac{c_0\alpha a n_0}{p}\right)\right\}
+\frac{\tmix}{\sqrt m} \notag\\
&\quad+\tmix^2\sqrt m\,\alpha^2
+(1+d^{1/q})p^{7/2}\tmix^{5/2}
\sqrt m\,\alpha^{3/2}L_\alpha^{3/2}
+\tmix p^{3/2}\sqrt\alpha \notag\\
&\quad+\frac{\tmix^{3/2}p^3L_\alpha^{1/2}}{\sqrt{\alpha m}}
+\frac{p^2\tmix^{3/2}aL_\alpha^{1/2}}{\sqrt m}
+\frac{pA_{\mathrm{st}}(p,q,\alpha)}
{\alpha a\sqrt m}
\exp\!\left(-\frac{c_{\rm st}\alpha a n_0}{p}\right).
\label{eq:BR}
\end{align}
Постоянные $c_0,c_{\rm st}>0$ не зависят от
$\tmix,n,p,\alpha$.
\end{lemma}

\begin{proof}
В \eqref{eq:composite-remainder} начальный продукт раскладываем как
$\Gamma_{1:k}^{(w)}=B_w^k+(\Gamma_{1:k}^{(w)}-B_w^k)$.
Геометрическая сумма детерминированной части даёт первое слагаемое в скобках
\eqref{eq:BR}, а \eqref{eq:product-stability} --- второе.
Пуассоновская граница контролируется величиной $\tmix/\sqrt m$.

Для стационарной версии третьего члена
\eqref{eq:composite-remainder} сначала рассматриваем $J^{(1,w)}$.
Обозначим через $\Pi_{J,w}$ совместный стационарный закон соответствующих
$J$-координат и цепи и положим
\[
\overline\psi_w(j,z)=\widetilde A(z)j
-\E_{\Pi_{J,w}}[\widetilde A(Z_1)J_0^{(0,w)}].
\]
Суммирование рекурсии по стационарному окну длины $m$ даёт
\begin{align*}
\sum_{k=0}^{m-1}\{J_k^{(1,w)}-\E_{\Pi_{J,w}}J_0^{(1,w)}\}
&=-\Abar^{-1}\sum_{k=1}^{m}
\overline\psi_w(J_{k-1}^{(0,w)},Z_k)\\[1mm]
&\quad+\frac1w\Abar^{-1}\{J_0^{(1,w)}-J_m^{(1,w)}\}.
\end{align*}
Proposition~2 из \cite{Levin2026} сокращает линейную по шагу часть средних
двух ветвей и оставляет $\tmix^2\sqrt m\,\alpha^2$.
Corollary~6 контролирует первую сумму в последнем тождестве и даёт
\[
\tmix p^{3/2}\sqrt\alpha
+\frac{\tmix^{3/2}p^3L_\alpha^{1/2}}{\sqrt{\alpha m}}.
\]
Для границы используем равномерную оценку
\[
\norm{J_k^{(1,w)}}_{L_p}
\lesssim p^2\tmix^{3/2}wa\,\log^{1/2}\!\frac1{wa}.
\]
\begingroup
\emergencystretch=3em
Она является упрощённой формой \cite[Lemma~8]{Levin2026} и восходит к
\cite[Proposition~10]{Durmus2025}. Лемма сформулирована для конечного запуска
из нуля; к стационарному закону переходим конечнопрошлым построением и леммой
Фату. После деления на $w\sqrt m$ эта оценка даёт шестое слагаемое
\eqref{eq:BR}. Propositions~8--9 из \cite{Levin2026} и неравенство Минковского
дают член второго порядка в \eqref{eq:BR}. Наконец, лемма~\ref{lem:startup}
переносит стационарную оценку на конечный старт. Это доказывает
\eqref{eq:remainder-bound}; из опубликованных работ импортируются оценки
отдельных компонент, а перенос после разогрева доказан локально.
\par
\endgroup
\end{proof}

\subsection{Основной результат}

Нам понадобится следующее элементарное неравенство.

\begin{lemma}[Сглаживание]\label{lem:smoothing}
Для любых вещественных случайных величин $X,Y$ и $p\geq1$
\begin{equation}\label{eq:smoothing}
\dK(X+Y,\mathcal N(0,1))
\leq \dK(X,\mathcal N(0,1))
+\frac{e}{\sqrt{2\pi}}\norm Y_{L_p}+e^{-p}.
\end{equation}
\end{lemma}

\begin{proof}
Для $t>0$ включение
$\{X+Y\leq x,\ |Y|\leq t\}\subseteq\{X\leq x+t\}$,
оценка $\sup_x|\Phi(x+t)-\Phi(x)|\leq t/\sqrt{2\pi}$ и неравенство
Маркова дают добавку $t/\sqrt{2\pi}+\E|Y|^p/t^p$.
Берём $t=e\norm Y_{L_p}$ и симметрично оцениваем второй хвост.
\end{proof}

Для краткости обозначим
\begin{equation}\label{eq:BM}
\begin{aligned}
\mathcal B_M(m,n,p;u)
 &=C_1(u,\tmix)
 p^{\frac{3p+1}{2(2p+1)}}m^{-\frac{p}{2(2p+1)}}\\
 &\quad+C_2(u,\tmix)\left\{
 \frac{\log(2n)}{\sqrt m}+p\,m^{-\frac{p}{2p+1}}\right\}.
\end{aligned}
\end{equation}

\begin{lemma}[Мартингальная аппроксимация]\label{lem:martingale-be}
Пусть $m\geq n/2$, $p\geq2$, $\sigma^2(u)>0$ и
\begin{equation}\label{eq:variance-lower}
m\alpha a\geq
\frac{2C_{\rm var}\norm u^2}{\sigma^2(u)}.
\end{equation}
Тогда
\[
\dK\!\left(
\frac{u^\top M_{n,n_0}^{\mathrm{bRR}}}
{\sqrt m\,\sigma_{n,n_0}(u)},\mathcal N(0,1)\right)
\lesssim\mathcal B_M(m,n,p;u).
\]
\end{lemma}

\begin{proof}
Приращения ограничены величиной
$\kappa(u,\tmix)\lesssim C_Q\norm u\,\tmix\norm\eps_\infty$.
Из \eqref{eq:variance-bound} и \eqref{eq:variance-lower} следует
$m\sigma_{n,n_0}^2(u)\geq m\sigma^2(u)/2$.
Применяем общую мартингальную оценку Берри--Эссеена
\cite[Lemma~21]{Samsonov2025} с детерминированным масштабом
$\sqrt m\,\sigma_{n,n_0}(u)$ и используем
лемму~\ref{lem:bracket}. Её три слагаемых имеют соответственно порядки
\[
\frac{\log(2n)}{\sqrt m},\qquad
p^{\frac{3p+1}{2(2p+1)}}m^{-\frac{p}{2(2p+1)}},\qquad
p m^{-\frac{p}{2p+1}},
\]
что даёт \eqref{eq:BM}.
\end{proof}

\begin{theorem}[Гауссовская аппроксимация для RR+PR]
\label{thm:gaussian}
Пусть выполнены условия 1--4 и $m\geq n/2$. Предположим, что
$\sigma^2(u)>0$ и $q\geq4p\geq8$,
$2\alpha\leq\alpha_{\mathrm{adm}}(p,q,\tmix)$ и \eqref{eq:variance-lower}.
Тогда
\begin{align}
&\dK\!\left(
\frac{\sqrt n\,u^\top
(\bar\theta_{n,n_0}^{\RR,\alpha}-\theta^\star)}
{\sigma(u)},\mathcal N(0,1)\right) \notag\\
&\qquad\lesssim
\mathcal B_M(m,n,p;u)
+\frac{\norm u}{\sigma(u)}
\mathcal B_R(m,n_0,p,q,\alpha)
+e^{-p}
+\frac{C_{\rm var}\norm u^2}{m\alpha a\,\sigma^2(u)}
+\frac{n_0}{n}.
\label{eq:master-BE}
\end{align}
Оценка равномерна по начальному распределению $Z_0$; скрытая константа не
зависит от $\tmix,n,p,\alpha$.
Для статистики с естественной оконной нормировкой $\sqrt m$ та же правая
часть справедлива без последнего слагаемого $n_0/n$.
\end{theorem}

\begin{proof}
Разделив \eqref{eq:full-decomp} на
$\sigma_{n,n_0}(u)$, получаем $X+Y$, где
\[
X=-\frac{u^\top M_{n,n_0}^{\mathrm{bRR}}}
{\sqrt m\,\sigma_{n,n_0}(u)},\qquad
Y=\frac{\mathcal R_{n,n_0}^{\RR}(u)}
{\sigma_{n,n_0}(u)} .
\]
Лемма~\ref{lem:martingale-be}, применённая к приращениям с обратным знаком,
даёт
\[
\dK(X,\mathcal N(0,1))\lesssim\mathcal B_M.
\]
По лемме~\ref{lem:smoothing} и
\eqref{eq:remainder-bound} добавка от $Y$ не превосходит второго и третьего
членов \eqref{eq:master-BE}. Условие на $m\alpha a$ и
лемма~\ref{lem:variance} дают
$\sigma_{n,n_0}(u)\geq\sigma(u)/\sqrt2$ и позволяют заменить
$\sigma_{n,n_0}(u)$ на $\sigma(u)$ с ценой четвёртого члена. Наконец,
переход от $\sqrt m$ к $\sqrt n$ умножает статистику на
$\sqrt{n/m}$; поскольку $m\geq n/2$,
$|\sqrt{n/m}-1|\lesssim n_0/n$. Стандартное сравнение
$\Phi(x/s)$ и $\Phi(x)$ завершает доказательство.
\end{proof}

\begin{corollary}[Сбалансированный режим]\label{cor:balanced}
Пусть выполнены условия 1--4 и фиксировано направление $u$ с
$\sigma^2(u)>0$.
Пусть $\alpha=c/\sqrt n$,
$p=\max\{2,\lceil\log n\rceil\}$,
$q=\max\{4p,\lceil\log(ed)\rceil\}$. Пусть $0<C_-\leq C_0$, причём
$C_-$ достаточно велика в зависимости от $c_0$ и $c_{\rm st}$.
Предположим, что
\begin{equation}\label{eq:burn-window}
C_-\frac{\log^2n}{\alpha a}
\leq n_0\leq
C_0\frac{\log^2n}{\alpha a},\qquad m\geq n/2.
\end{equation}
Если $n$ достаточно велико и
$2\alpha\leq\alpha_{\mathrm{adm}}(p,q,\tmix)$, то
\begin{equation}\label{eq:balanced-BE}
\begin{aligned}
&\dK\!\left(
 \frac{\sqrt n\,u^\top
 (\bar\theta_{n,n_0}^{\RR,\alpha}-\theta^\star)}
 {\sigma(u)},\mathcal N(0,1)\right)\\
&\qquad\leq C(u,c,\theta_0,C_0,C_-)\,
 \frac{C_{\mathrm{mix}}(\tmix)\operatorname{polylog}(n)}
 {n^{1/4}} .
\end{aligned}
\end{equation}
Здесь $C_{\mathrm{mix}}(\tmix)<\infty$ собирает зависимость констант
мартингальной оценки и моментных неравенств от времени смешивания; её
полиномиальность не предполагается.
\end{corollary}

\begin{proof}
Подставим $\alpha,p,q$ в
\eqref{eq:BM}--\eqref{eq:master-BE}. Для достаточно большого $n$ имеем
$m\alpha a\asymp\sqrt n$, поэтому условие~\eqref{eq:variance-lower}
выполнено автоматически. Условие~\eqref{eq:burn-window} делает оба переходных
экспоненциальных члена не выше полилогарифмического множителя при
$n^{-1/4}$. Остальные слагаемые имеют порядок не больше
$C_{\mathrm{mix}}(\tmix)\operatorname{polylog}(n)n^{-1/4}$;
цены замены дисперсии и нормировки равны соответственно $O(n^{-1/2})$ и
$O(n^{-1/2}\log^2n)$.
\end{proof}

\begin{remark}[Почему нужна общая траектория]\label{rem:shared}
Внутри PR-окна веса обеих ветвей близки к $\Abar^{-1}$. На общей траектории
один и тот же шум входит с коэффициентом $2-1=1$, поэтому сохраняется
$\Sigma_\infty$. Для независимых траекторий ведущие ковариации сложились бы
как $2^2\Sigma_\infty+(-1)^2\Sigma_\infty=5\Sigma_\infty$.
\end{remark}

